% ============================================================================
% Fichier  : partie2/P02_C11_traiter.tex
% Titre    : Traiter, Suivre et Rendre Compte
% Auteur   : MINKA MI NGUIDJOÏ Thierry Emmanuel
% Version  : 1.0 -- novembre 2026
% ============================================================================
% Sources : ISO 31000:2018 (6.5 à 6.7) ; COSO ERM (2017) ; ISO/IEC 27005.
% Plan de traitement de LiZbiZ : cas fictif, montants cohérents entre eux.
% ============================================================================

\chapter{Traiter, Suivre et Rendre Compte}
\label{chap:traiter}
\chaptermark{Traiter, suivre, rendre compte}

\epigraph{%
    « Un risque n'est traité\\
    que lorsque la mesure fonctionne. »%
}{Principe du cours}

\niveauChapitre{Intermédiaire}{Chapitres~\ref{chap:gouvernance},
\ref{chap:matrices} et~\ref{chap:quantitatif}.}

\begin{boxobjectifs}
\begin{itemize}
    \item Choisir une option de traitement selon l'ISO~31000.
    \item Rédiger un plan de traitement complet.
    \item Arbitrer un budget par les pertes évitées.
    \item Définir des indicateurs clés de risque avec des seuils.
    \item Rendre compte au comité et au conseil.
\end{itemize}
\end{boxobjectifs}

\bigskip

% ============================================================================
\section*{Question directrice : comment dépenser les 180 millions votés en avril~?}
% ============================================================================

En avril, sous le choc, le comité de direction a voté 180~millions de
FCFA pour la sécurité. En juin, le registre et les estimations sont
prêts. Il reste à choisir les traitements, à les confier à des
responsables, et à savoir, chaque trimestre, si le risque baisse
vraiment.

\medskip
\textbf{Plan du chapitre.}
\begin{enumerate}
    \item Les options de traitement (section~\ref{sec:c11-options}).
    \item Le plan de traitement (section~\ref{sec:c11-plan}).
    \item Arbitrer le budget (section~\ref{sec:c11-budget}).
    \item Les indicateurs clés de risque (section~\ref{sec:c11-kri}).
    \item Rendre compte (section~\ref{sec:c11-rendre}).
\end{enumerate}

\separateur

% ============================================================================
\section{Les Options de Traitement}
\label{sec:c11-options}
% ============================================================================

L'ISO~31000 énumère plusieurs options, qui se combinent~\cite{iso31000_2018}.

\begin{tableaucours}{L{3.6cm}L{4.6cm}L{4.0cm}}{%
    \tableauHeader{Option} &
    \tableauHeader{Principe} &
    \tableauHeader{Exemple chez LiZbiZ}
}
Éviter & Renoncer à l'activité qui crée le risque & Supprimer l'accès
distant permanent du prestataire \\
Retirer la source & Supprimer la cause & Remplacer FTP par SFTP \\
Modifier la vraisemblance & Rendre l'événement moins probable & Second
facteur, correctifs \\
Modifier les conséquences & Limiter les dommages & Sauvegardes hors ligne,
mode dégradé \\
Partager & Transférer une part des conséquences & Assurance cyber, clauses
du contrat \\
Maintenir & Accepter en connaissance de cause, par écrit & Défiguration du
portail, sous surveillance \\
Prendre le risque & Saisir une opportunité & Ouvrir un portail client en
ligne, avec ses mesures \\
\end{tableaucours}
\vspace{-0.8em}
\begin{center}{\small\itshape Tableau~11.1. Les options de traitement (d'après l'ISO 31000:2018)}\end{center}

\begin{boxalert}
\textbf{Violation typique.} Croire que l'assurance « traite » le risque.

\smallskip
\textbf{Règle~:} l'assurance partage une part des pertes financières~;
elle ne réduit ni la vraisemblance, ni l'arrêt de l'activité, ni l'atteinte
aux clients. Elle complète la réduction, elle ne la remplace pas. Et
l'assureur exige souvent, en contrepartie, des mesures (second facteur,
sauvegardes).
\end{boxalert}

% ============================================================================
\section{Le Plan de Traitement}
\label{sec:c11-plan}
% ============================================================================

\begin{boxdef}
\termecle{Plan de traitement}~: document qui précise, pour chaque risque
traité, les actions retenues, leur justification, le responsable, les
ressources, l'échéance, l'indicateur de suivi et le risque résiduel visé.
\end{boxdef}

\begin{tableaucours}{L{1.0cm}L{4.2cm}L{2.2cm}L{1.6cm}L{2.8cm}}{%
    \tableauHeader{Risque} &
    \tableauHeader{Action} &
    \tableauHeader{Responsable} &
    \tableauHeader{Échéance} &
    \tableauHeader{Indicateur}
}
R-02 & Second facteur sur la passerelle & M.~KOUNOU & Juillet & Part des
connexions avec second facteur \\
R-02 & Bastion d'administration & M.~KOUNOU & Septembre & Part des accès
du prestataire par le bastion \\
R-01 & Sauvegardes hors ligne et immuables & M.~KOUNOU & Août & Durée du
dernier test de restauration \\
R-07 & Chiffrement des bases clients & M.~KOUNOU & Octobre & Part des
bases chiffrées \\
R-11 & Rayonnages surélevés et pompes & M.~NGUELE & Septembre & Hauteur
des stocks au sol \\
R-03 & Second lien télécom à Garoua & M.~NGUELE & Novembre & Durée des
coupures \\
\end{tableaucours}
\vspace{-0.8em}
\begin{center}{\small\itshape Tableau~11.2. Extrait du plan de traitement de LiZbiZ}\end{center}

Le \termecle{risque résiduel} est réévalué avec les seules mesures en place
et vérifiées~; il est \emph{accepté par écrit} par le propriétaire et, au
besoin, par la direction. Une mesure décidée mais non mise en œuvre ne
réduit rien (cours SEC3266, Chapitre sur les risques).

% ============================================================================
\section{Arbitrer le Budget}
\label{sec:c11-budget}
% ============================================================================

On classe les traitements par \emph{pertes annuelles évitées nettes} et par
rapport au coût annuel. Les pertes évitées viennent des estimations FAIR
(Chapitre~\ref{chap:quantitatif}) pour les risques majeurs, et d'ordres de
grandeur pour les autres.

\begin{tableaucours}{L{4.8cm}R{1.8cm}R{2.2cm}R{2.4cm}}{%
    \tableauHeader{Traitement} &
    \tableauHeader{Coût annuel} &
    \tableauHeader{Pertes évitées} &
    \tableauHeader{Rapport net}
}
Second facteur (R-02) & 6 & 74 & 11,3 \\
Sauvegardes hors ligne (R-01) & 15 & 60 & 3,0 \\
Bastion (R-02) & 10 & 25 & 1,5 \\
Rayonnages et pompes (R-11) & 4 & 9 & 1,3 \\
Surveillance externalisée (SOC) & 60 & 45 & -0,25 \\
\end{tableaucours}
\vspace{-0.8em}
\begin{center}{\small\itshape Tableau~11.3. Arbitrage des traitements (millions de FCFA par an ; rapport = (évité moins coût) / coût)}\end{center}

\begin{boxscenario}{Opération LIZBIZ. La décision de juin}
\textbf{Constatation.} Les quatre premiers traitements coûtent 35~millions
par an et évitent environ 168~millions de pertes attendues. La
surveillance externalisée coûte plus qu'elle n'évite, selon les
estimations.

\medskip
\textbf{Analyse} (confiance faible). La surveillance réduit surtout le
délai de détection, dont l'effet sur les pertes est mal estimé~; le
rapport négatif peut venir de cette sous-estimation. Le comité retient les
quatre premiers traitements, et une surveillance limitée aux heures de
nuit, à réévaluer après six mois de mesures.
\end{boxscenario}

% ============================================================================
\section{Les Indicateurs Clés de Risque}
\label{sec:c11-kri}
% ============================================================================

\begin{boxdef}
\termecle{Indicateur clé de risque} (KRI)~: mesure qui signale à l'avance
l'évolution d'un risque, assortie de seuils d'alerte. Il se distingue de
l'indicateur de performance, qui mesure l'atteinte d'un objectif.
\end{boxdef}

\begin{tableaucours}{L{4.8cm}L{1.0cm}L{2.0cm}L{2.0cm}L{2.0cm}}{%
    \tableauHeader{Indicateur} &
    \tableauHeader{Risque} &
    \tableauHeader{Vert} &
    \tableauHeader{Orange} &
    \tableauHeader{Rouge}
}
Connexions distantes avec second facteur & R-02 & 100~\% & 95 à 99~\% &
moins de 95~\% \\
Âge du dernier test de restauration réussi & R-01 & moins de 3~mois & 3 à
6~mois & plus de 6~mois \\
Correctifs critiques en retard & R-02 & 0 & 1 à 3 & plus de 3 \\
Comptes orphelins actifs & R-07 & 0 & 1 à 5 & plus de 5 \\
\end{tableaucours}
\vspace{-0.8em}
\begin{center}{\small\itshape Tableau~11.4. Indicateurs clés de risque de LiZbiZ}\end{center}

Un bon indicateur clé de risque est \emph{avancé} (il bouge avant
l'incident), \emph{mesurable} automatiquement si possible, et rattaché à un
\emph{propriétaire} qui agit quand il passe à l'orange.

% ============================================================================
\section{Rendre Compte}
\label{sec:c11-rendre}
% ============================================================================

\begin{boxprocedure}
\textbf{Le rapport trimestriel au comité des risques}
\begin{enumerate}
    \item Les dix premiers risques, avec leur score et leur évolution
          depuis le trimestre précédent.
    \item L'avancement du plan de traitement~: actions faites, en retard,
          bloquées.
    \item Les indicateurs clés de risque passés à l'orange ou au rouge, et
          ce qui a été fait.
    \item Les incidents du trimestre et les risques qu'ils révèlent.
    \item Les risques nouveaux ou émergents.
    \item Les décisions demandées au comité.
\end{enumerate}
\end{boxprocedure}

\begin{boxobjet}{R-01 · Arrêt de la facturation (suivi)}
\textbf{Cotation~:} V=3, I=3 (score~9) en juin~; cible V=2, I=2 après
sauvegardes testées et mode dégradé. \textbf{Traitement~:} sauvegardes hors
ligne (août), test trimestriel. \textbf{Propriétaire~:} Mme~MBALLA.
\textbf{Indicateur~:} âge du dernier test de restauration réussi (vert sous
trois mois).
\end{boxobjet}

\begin{boxmetier}{Gestionnaire de risques}
\textbf{Une page pour le conseil.} Le conseil reçoit chaque semestre une
seule page~: respect de l'appétence, les cinq risques majeurs, les
décisions qui lui reviennent. Le détail est en annexe.
\end{boxmetier}

\separateur

\begin{boxresume}
\textbf{Ce qu'il faut retenir de ce chapitre~:}
\begin{itemize}
    \item \textbf{Plusieurs options} se combinent~; l'assurance partage,
          elle ne réduit pas.
    \item \textbf{Le plan de traitement} nomme l'action, le responsable,
          l'échéance, l'indicateur et le résiduel visé.
    \item \textbf{On arbitre} par les pertes évitées nettes, en disant la
          confiance de l'estimation.
    \item \textbf{Les indicateurs clés de risque} sont avancés, mesurables,
          à seuils, et ont un propriétaire.
    \item \textbf{Rendre compte}~: dix risques, plan, indicateurs,
          incidents, émergents, décisions.
\end{itemize}
\end{boxresume}

\begin{boxinfo}
\textbf{Points de contrôle de l'élève.} Je sais~: choisir une option de
traitement~; rédiger une ligne de plan de traitement~; calculer un rapport
net~; définir un indicateur clé de risque et ses seuils~; structurer un
rapport trimestriel.
\end{boxinfo}

\bigskip

\begin{boxexo}[title={\ding{45}\quad Exercice~11.1. Quelle option~? \niveaubase}]
Donnez l'option de traitement~: souscrire une assurance~; fermer l'agence
de Garoua~; installer des pompes à Bonabéri~; accepter la défiguration du
portail.
\end{boxexo}

\begin{boxexo}[title={\ding{45}\quad Exercice~11.2. Rapport net \niveaubase}]
Un traitement coûte 8~millions par an et évite 20~millions de pertes
attendues. Calculez le rapport net.
\end{boxexo}

\begin{boxexo}[title={\ding{45}\quad Exercice~11.3. KRI ou KPI \niveaubase}]
Classez en indicateur clé de risque ou de performance~: le chiffre
d'affaires mensuel~; le nombre de correctifs critiques en retard~; le taux
de livraison à l'heure.
\end{boxexo}

\begin{boxexo}[title={\ding{45}\quad Exercice~11.4. L'assurance \niveaumoyen}]
Pourquoi l'assurance seule ne suffit-elle pas pour R-02~? Que demandera
probablement l'assureur à LiZbiZ~?
\end{boxexo}

\begin{boxexo}[title={\ding{45}\quad Exercice~11.5. Une ligne de plan \niveaumoyen}]
Rédigez la ligne du plan de traitement pour le chiffrement des
ordinateurs portables de la direction.
\end{boxexo}

\begin{boxexo}[title={\ding{45}\quad Exercice~11.6. Le budget \niveaumoyen}]
Sur les 180~millions votés, combien restent après les quatre premiers
traitements du Tableau~11.3 la première année~? Que proposez-vous d'en
faire~?
\end{boxexo}

\begin{boxexo}[title={\ding{45}\quad Exercice~11.7. Le SOC \niveaumoyen}]
Le rapport net de la surveillance externalisée est négatif. Quelles
pertes évitées l'estimation a-t-elle pu oublier~?
\end{boxexo}

\begin{boxexo}[title={\ding{45}\quad Exercice~11.8. Seuils \niveaumoyen}]
Proposez des seuils vert, orange et rouge pour l'indicateur « hauteur des
marchandises stockées au sol à Bonabéri en saison des pluies ».
\end{boxexo}

\begin{boxexo}[title={\ding{45}\quad Exercice~11.9. Le rapport \niveauexpert}]
Rédigez le rapport trimestriel de septembre au comité des risques de
LiZbiZ, à partir des éléments du chapitre.
\end{boxexo}

\begin{boxexo}[title={\ding{45}\quad Exercice~11.10. La page du conseil \niveauexpert}]
Rédigez la page semestrielle au conseil~: appétence, cinq risques majeurs,
décisions demandées.
\end{boxexo}

\bigskip

\begin{boxinfo}
\textbf{Pour aller plus loin.}
\begin{itemize}
    \item ISO~31000, traitement, suivi et rapport~\cite{iso31000_2018}.
    \item COSO ERM~\cite{coso2017erm}.
\end{itemize}
\end{boxinfo}

\begin{boxinfo}
\textbf{Transition.} Un dispositif de risques se vérifie. La Partie~III
apprend à l'auditer, en commençant par l'ISO~19011
(Chapitre~\ref{chap:audit}).
\end{boxinfo}
